În zilele noastre, transportul public este rar utilizat din cauza lipsei de confort, a ineficienței sau a neîndeplinirii așteptărilor utilizatorilor. Problemele frecvente includ autobuze care nu respectă orarele sau dificultăți în achiziționarea biletelor. În multe locuri, orarele nu sunt afișate în stații și nici disponibile online. Se întâmplă adesea ca autobuzele să fie aglomerate, iar călătorii află acest lucru doar când ajung în stație, și nu au alte opțiuni de transport. Platforma online Bus4U își propune să rezolve aceste probleme.

Bus4U constă dintr-o aplicație mobilă, un site web pentru utilizatori și o interfață web de administrare. Funcțiile sale principale includ căutarea rutelor de autobuz prin specificarea punctului de plecare, a destinației și a orei, achiziționarea biletelor online în avans și vizualizarea orelor de plecare defalcate pe orașe și stații. În plus, stațiile de autobuz pot fi vizualizate pe o hartă filtrată după orașe și rute, iar pozițiile actuale ale autobuzelor pot fi urmărite în timp real. Pentru a asigura aceste funcționalități și gestionarea biletelor, o aplicație dedicată rulează pe terminalele POS instalate în autobuze. Această aplicație validează biletele și partajează informațiile live despre autobuze cu serverul central.

În plus, aplicația web oferă o interfață de administrare pentru companiile de autobuze. Aici, administratorii pot gestiona orarele, stațiile și încărca rutele. Administratorii companiilor pot gestiona, de asemenea, informațiile angajaților și diverse date despre autobuze, cum ar fi inspecțiile tehnice periodice și asigurările. Este disponibilă și o interfață statistică, care permite conducătorilor companiilor să vizualizeze numărul de utilizatori înregistrați, biletele vândute și veniturile generate, defalcate pe lună, an sau pe întreaga perioadă.

Posibilitățile de dezvoltare viitoare includ introducerea abonamentelor și a diverselor reduceri, precum și urmărirea în timp real a gradului de ocupare a autobuzelor. Din perspectiva companiilor de autobuze, monitorizarea programului de lucru al șoferilor și implementarea contabilității automate ar putea fi, de asemenea, benefice.

Această lucrare se concentrează pe interfața utilizatorului a software-ului, incluzând cele două site-uri web și aplicația mobilă, în timp ce funcționalitatea backend este acoperită într-o altă lucrare.

\textbf{Cuvinte cheie:} transport în comun, cursă de autobuz, digitalizare